\section{Compact propagation with possible Jordan blocks}\label{c7:compact}

The horizontal gas gives a representation of the cylinder partition, but
its weights are signed. For the balanced lists with $m$ sites in each
group, we now turn that representation into a compact operator and
estimate the effect of the site factors on its power traces.
The main point is that these factors suppress particles on an interval of
length proportional to $\log m$. We compare propagation across that
interval with propagation outside it, while retaining every peripheral
Jordan term of the exterior operator.

There are two inputs from the preceding sections. Absolute stability
controls particle counts and the common charge on each spatial block.
The residual pair interaction has an exponentially convergent expansion
at large separation. The first input gives a summable local kernel after
a change of gauge; the second is carried exactly by an auxiliary
Bargmann space. We first construct the state space and establish the
bounds used by the operator representation.

\subsection{The local state and its gauge}

Fix a sufficiently small positive post-collision displacement $d$, as
in Section~\ref{c7:stability}. Take $T=bn$, where the integer block size
$b$ will be chosen sufficiently large and then held fixed. In block $j$
write $x=jb+t$, $0<t<b$. Positions at ties or block boundaries form null
sets and are discarded.

A block configuration $Y$ is a finite ordered list of positions $t$ and
particle labels. It is measured by Lebesgue measure on each ordered
simplex and counting measure on labels; the empty configuration is an
atom. This is equivalent to the factorial convention for indistinguishable
particles. Write $q(Y)$ for its total charge vector. The state before a
block is
\[
 S=(Y_{\rm prev},p),\qquad
 p\in\mathbb Z^2/8\mathbb Z(1,-1),
\]
where $Y_{\rm prev}$ records the preceding block in translated
coordinates. Its successor after inserting $Y$ is
\begin{equation}\label{j:successor-state}
 S'=(Y,p-q(Y)).
\end{equation}
We use the indicated charge jumps to choose one consistent lift within
each block. State measure is the product of the block-configuration
measure and counting measure on $p$. Particle types with zero local
factor may be omitted. In particular $P,M$ in each species have nonzero
local factors, so their charges can be used to connect states.

Separate the forward residual interactions into pairs in the same block,
pairs in neighboring blocks, and pairs at least two blocks apart. Let
$w_\theta(S,S')$ be the product of all constants, phases, and the
integral-square weight in $Y$, starting at $p$, together with residual
factors within $Y$ and from $Y_{\rm prev}$ to $Y$. The site factors
involving $B_T$ are omitted for now. The function $w_\theta$ is continuous
and nonzero on each allowed open position simplex.

For a closed sequence of $n$ blocks, the sum of the absolute values of
the omitted tail log-magnitudes is at most
\[
 C e^{-cb}\sum_j |Y_j|^2.
\]
Indeed an interaction across $l\ge2$ blocks has horizontal separation at
least $(l-1)b$, and each residual log decays exponentially there. This
also applies to short periodic sequences, with their images included.
For large $b$, stability therefore bounds the sum of the local costs by
\begin{equation}\label{j:local-cycle-cost}
 \sum_j-\log|w_0(S_j,S_{j+1})|
 \ge \frac{c'}b\sum_j\bigl(|Y_j|^2+p_{c,j}^2\bigr)-C_b n,
 \qquad p_{c,j}=\frac{p_{1,j}+p_{2,j}}{\sqrt2}.
\end{equation}
Fix an open parameter neighborhood $\mathcal U_\theta$ whose closure $D$ is compact in
the complex $\theta$ plane.
The same bound, with adjusted constants, holds for the sum of
$-\log\sup_{\theta\in D}|w_\theta|$. The change from $\theta=0$ costs at
most a constant times $b(1+|p_c|+|Y|)$ per block, which is absorbed by
the quadratic terms. Constants in this section may depend on $b,d,D$.

There is consequently a real measurable function $\Phi$ on states such
that
\begin{equation}\label{j:local-gauge}
 \left|e^{\Phi(S')-\Phi(S)}w_\theta(S,S')\right|
 \le \exp\!\left(C_b-\delta_b
       (|Y_{\rm prev}|^2+|Y|^2+p_c^2)\right),
 \qquad \delta_b=c''/b>0,
\end{equation}
for $\theta\in D$ and allowed successors. To construct it, subtract a
smaller quadratic cost from the edge costs used in
\eqref{j:local-cycle-cost}, and add $C_b$, so that every finite directed
cycle has nonnegative modified cost. Let $\Phi(S)$ be the infimum of
modified costs of paths from a reference state to $S$. Every state can
be reached and returned from: insert $P,M$ charges as needed at distinct
positions, and then insert the required final block. A fixed return path
bounds the infimum below by the negative of its cost, so the infimum is
finite. The directed triangle inequality gives \eqref{j:local-gauge}.
For measurability, the intermediate positions in the infimum may be
restricted to rational translated coordinates, using continuity on the
open simplices; the terminal state $S$ remains arbitrary and fixed.
The discrete constraints depend only on the labels. On each label
stratum this is a countable infimum of functions continuous in the
endpoint, hence is measurable.

\subsection{Exact propagation of the interaction tail}

Choose the logarithm of each far residual factor using the expansion of
$K^r$ and the convergent logarithmic expansions of the short and outer
factors. For an earlier label $u$, a later label $v$, and sufficiently
large positive separation $s$, that log has a representation
\begin{equation}\label{j:tail-linear-form}
 g_v e^{-D_0s}h_u
\end{equation}
on a complex sequence space, with bilinear row-column pairing. Here
$D_0$ is a direct sum of finite Jordan arrays whose diagonal values are
$\pi l/6$, $l\ge1$. Their sizes and the entries of their nilpotent parts
are uniformly bounded; those parts generate the polynomial factors below. The components
of $g_v,h_u$ have at most polynomial growth.
To obtain this representation, write the residual log as
\[
 \sum_{l\ge1}e^{-\pi l s/6}P_{l;uv}(s).
\]
The polynomials have uniformly bounded degree and polynomially growing
coefficients, also after the fixed imaginary height shifts. A nilpotent
chain generates each power of $s$; separate coordinates can be used for
each ordered label pair and degree. Any positive intervening delay then
puts the resulting vectors in $\ell^2$.

Set
\begin{equation}\label{j:tail-data}
 \begin{split}
 A_0&=e^{-bD_0},\\
 G(Y)&=\sum_{(t,v)\in Y}g_v e^{-D_0(t+b/2)},\\
 H_1(Y_{\rm prev})&=\sum_{(t,u)\in Y_{\rm prev}}
                    e^{-D_0(3b/2-t)}h_u.
 \end{split}
\end{equation}
Thus $\|G(Y)\|\le Ce^{-cb}|Y|$ and
$\|H_1(Y_{\rm prev})\|\le Ce^{-cb}|Y_{\rm prev}|$, while the trace norm
of $A_0$ tends to zero as $b$ increases. The offsets in
\eqref{j:tail-data} put exactly two blocks between a source and its first
use. For example, a particle $(t_u,u)$ in $Y_{j-2}$ enters the field used
by $(t_v,v)$ in $Y_j$ with exponent
\[
 g_v e^{-D_0(t_v+b/2)}e^{-D_0(3b/2-t_u)}h_u
 =g_v e^{-D_0(2b+t_v-t_u)}h_u,
\]
which is \eqref{j:tail-linear-form} at their actual separation. Each
additional factor $A_0$ adds one more block of delay.

On the holomorphic Bargmann space over this sequence space, with
orthonormal monomials $z^\alpha/\sqrt{\alpha!}$, define the edge operator
\begin{equation}\label{j:bargmann-edge}
 \mathcal K(S,S')f(z)=e^{G(Y)z}f(A_0z+H_1(Y_{\rm prev})).
\end{equation}
Finite-mode Bose and Fermi trace--determinant formulas are derived in
\cite[Section~2]{Klich2002}. Here the affine terms and infinite-mode limit
are also needed.

For large $b$, the trace norm of \eqref{j:bargmann-edge} is at most
\begin{equation}\label{j:bargmann-bound}
 C\exp\!\left(C(\|G(Y)\|^2+\|H_1(Y_{\rm prev})\|^2)\right).
\end{equation}
Indeed linear composition $C_Bf(z)=f(Bz)$ acts by symmetric tensor
powers and has trace norm at most $(1-\|B\|_1)^{-1}$ when $\|B\|_1<1$.
Factor \eqref{j:bargmann-edge} through $C_{4A_0}$, with outside linear
dampings $z\mapsto z/2$ carrying the exponential multiplier and the
translation. Such damped operators are bounded with the estimate in
\eqref{j:bargmann-bound}: a map $e^{gz}f(dz+h)$, $d=1/2$, is a scalar
times $W_vC_dW_w$, where
\[
 W_vf(z)=e^{-\|v\|^2/2+\bar v^tz}f(z-v),
 \qquad v+dw=\bar g,\quad dv+w=-h.
\]
The operators $W_v$ are unitary and strongly continuous, by Gaussian
translation in finite dimensions or directly on total exponential
vectors.

For a closed sequence, multiply the edge operators in increasing block
order. Composition and trace give
\begin{equation}\label{j:bargmann-cycle}
 \begin{gathered}
 \operatorname{tr}\prod_{j=0}^{n-1}\mathcal K(S_j,S_{j+1})
 =\det(I-A_0^n)^{-1}\exp\!\left(\sum_jG(Y_j)z_j\right),\\
 z_{j+1}=A_0z_j+H_1(Y_{j-1}),\qquad z_n=z_0.
 \end{gathered}
\end{equation}
In finite one-particle dimension, the trace of $e^{gz}f(Bz+h)$ for
small $\|B\|$ is
$\det(I-B)^{-1}\exp(g(I-B)^{-1}h)$. Integrating its coherent diagonal
$e^{gz+\bar z^t(Bz+h)}$ against normalized complex Gaussian gives this
formula: Gaussian monomial orthogonality gives the trace resolution, and
square completion works first for Hermitian $B$ and then analytically in
its entries. Composition proves \eqref{j:bargmann-cycle} in finite
dimension. Increasing Jordan-block truncations converge in trace norm by
the preceding factorization: $C_{4A_0}$ converges in trace norm by its
tensor series, and the outside factors and their adjoints converge
strongly. The periodic solution $z_j$ sums the sources in blocks up to
$j-2$ in the periodic past. The explicit delay calculation above shows
that the exponential in \eqref{j:bargmann-cycle} is precisely the entire
omitted tail, including its images.

\begin{proposition}[Block kernel and trace realization]\label{j:block-kernel}
Use the state space, measure, and successor rule
\eqref{j:successor-state}. Suppose the local kernel satisfies
\eqref{j:local-gauge}, and use the tail data and edge operator
\eqref{j:tail-data}--\eqref{j:bargmann-edge} with the bounds just proved.
On $L^2$ of states with these Bargmann fibers, give an allowed edge the
operator-valued kernel
\begin{equation}\label{j:gauged-kernel}
 e^{\Phi(S')-\Phi(S)}w_\theta(S,S')\,\mu(Y)\,\mathcal K(S,S'),
\end{equation}
and give every other edge value zero. The choices of $\mu(Y)$ used in this
section are one, the indicator that every label is $M$ or $F_{\rm pair}$,
and the products of the specified balanced periodized or infinite-volume
$B$ factors. Their per-particle factors have modulus at most one. For
sufficiently large fixed $b$, these kernels are Hilbert--Schmidt uniformly
in those choices and in $\theta\in D$.
The same conclusion permits fixed polynomial factors in the particle
counts. For the fixed gauge on $D$, the kernels form holomorphic families
in Hilbert--Schmidt norm on the open parameter neighborhood $\mathcal U_\theta$, with
bounds uniform on compact subsets.

For a product of $n\ge2$ such kernels, its trace is the absolutely
convergent integral over closed successor sequences, with the gauge
factors canceled and the Bargmann trace given by
\eqref{j:bargmann-cycle}. In particular, let $T_j$ carry the specified
periodized $B_T$ site factors for block $j$. Then, in increasing block order,
\begin{equation}\label{j:operator-gas-trace}
 \operatorname{tr}\prod_{j=0}^{n-1}T_j
 =\det(I-A_0^n)^{-1}Z_{m,bn}.
\end{equation}
\end{proposition}
\begin{proof}
The squared Hilbert--Schmidt norm is the integral of the squared fiber
Hilbert--Schmidt norms. By \eqref{j:bargmann-bound}, their extra cost is
at most $\exp(Ce^{-2cb}(|Y_{\rm prev}|^2+|Y|^2))$. Increase $b$ so this
is absorbed by \eqref{j:local-gauge}. The remaining integrals converge:
the volume for a list of $r$ positions is $b^r/r!$, there are finitely
many labels, and the charge sum has finitely many representatives for
each value of $p_1+p_2$. The successor charge $p'$ is uniquely fixed by
$Y,p$, so no additional charge sum occurs in the squared kernel integral.
The same domination absorbs fixed polynomial count factors and gives the
stated local holomorphy.

For two kernels, the trace formula is Hilbert--Schmidt pairing. For a
longer finite product, the same confining bounds make every successive
contraction over states and fibers absolute. The gauges telescope on a
closed sequence. Equation~\eqref{j:bargmann-cycle} then supplies exactly
the far interactions, while $w_\theta$ supplies the local ones, proving
\eqref{j:operator-gas-trace}.
\end{proof}

The marked profile comparisons will apply this construction after
establishing their local bounds and suppression estimates.

\subsection{Empty-block traces and spectral consequences}

Let $E_0$ be the homogeneous operator with every site multiplier equal
to one. Let $I_0$ be the operator retaining only the labels
$M,F_{\rm pair}$. Both depend on $\theta$. The exact traces are
\begin{equation}\label{j:empty-block-traces}
 \begin{split}
 \operatorname{tr}E_0^n
   &=\det(I-A_0^n)^{-1}/c(bn),\\
 \operatorname{tr}I_0^n
   &=8\det(I-A_0^n)^{-1}\sum_{k\in\mathbb Z}u_k^n,
 \end{split}
 \qquad n\ge2.
\end{equation}
Here
\[
 u_k=\exp\!\left[-\frac{4\pi b}{9}(k+\theta/\pi)^2\right].
\]
The first equality follows from the empty contour at $m=0$ and the gas
identity. In the second, every retained particle has strictly negative
common charge. A closed sequence can therefore contain no particles,
and the eight representatives at each $k=p_1+p_2$ give the displayed
empty-block sum. The same reasoning gives the trace formula with the sum
restricted to a consecutive range of $k$ for a triangular diagonal
compression to that range. This full formula will also be used for the
marked intervals.

We recall the spectral facts needed to read these traces. For a compact
Hilbert--Schmidt operator, the limsup of the absolute $n$th roots of its
power traces is its spectral radius, and those traces determine every
nonzero eigenvalue with its algebraic multiplicity. To see this,
approximate the operator by finite rank with norm error less than
$\epsilon$. Outside radius $\epsilon$, the error resolvent exists and
finite-dimensional inversion makes the full resolvent meromorphic with
finite-rank pole parts. Resolvent contour integration isolates the Riesz
blocks. The resolvent identity makes each Riesz projection idempotent and
commuting, and Laurent comparison identifies the higher pole parts with
powers of the nilpotent part on that block. The complementary powers are
exponentially bounded at any radius containing the remaining spectrum;
their traces have the same rate by retaining two Hilbert--Schmidt factors.
On an isolated spectral radius the trace is therefore a sum of algebraic
multiplicities times eigenvalue powers. Ces\`aro extraction of distinct
unit-circle frequencies shows that this sum cannot tend to zero unless
all its coefficients vanish. Applying this successively off zero proves
the assertions. Separating resolvent contours give locally uniform bounds
for norm-continuous families.

It follows from \eqref{j:empty-block-traces} and
$\log c(T)/T\to\pi/16$ that all nonzero eigenvalues of $E_0$ are
independent of $\theta$, and
\begin{equation}\label{j:exterior-radius}
 r(E_0)=r_E=e^{-b\pi/16}.
\end{equation}
For $|\Re\theta|<\pi/2$, the unique leading eigenvalue of $I_0$ is
$u=u_0$, with algebraic multiplicity eight. It is semisimple in this
strip. Along successor edges a negative particle charge increases $k$;
the operator acts from an argument state to an output state, so in that
operator direction $k$ can only stay or decrease. Split into the three
triangular diagonal compressions with negative, zero, and positive $k$.
Their traces are \eqref{j:empty-block-traces} restricted to the
corresponding ranges, so the two side blocks are spectrally separated
from $u$. In the zero block, a successor state must have empty $Y$.
Splitting off these empty fibers leaves, per representative, the diagonal
operator $uC_{A_0}$ on an empty fiber and zero on the other diagonal block.
The constants give a simple semisimple eigenvalue $u$ for $uC_{A_0}$;
$C_{A_0}$ is strictly contracting on higher Bargmann degrees.

For completeness, semisimplicity passes through each of these finite
triangular decompositions as follows. A generalized $u$-eigenvector has
zero component on the incoming nonresonant side by invertibility there.
Its component on the resonant block is already a genuine eigenvector.
Applying the operator minus $u$ leaves a generalized eigenvector on the
outgoing nonresonant side, which must also be zero. Thus, locally uniformly
on compact neighborhoods in the strip,
\begin{equation}\label{j:interior-powers}
 (I_0/u)^s=P_I+O(e^{-cs})
\end{equation}
in operator norm, where $P_I$ is its Riesz projection. For
$E'=E_0/r_E$, peripheral Jordan blocks are retained: locally uniformly,
$\|(E')^s\|\le C(1+s)^r$, and its powers equal their finite peripheral
Riesz expansions up to exponentially small error. The uniformity follows
from fixed separating contours and the fixed finite ranks of the Riesz
blocks.

\subsection{The logarithmic interval and its exterior words}

We now return to the balanced $2m$-site problem, with $m$ even and
sufficiently large at fixed $b$. Put
\[
 \kappa_1=\pi/3,\qquad
 L=\left\lfloor\frac{\log m}{\kappa_1b}\right\rfloor,
 \qquad a_m=me^{-\kappa_1bL}\in[1,e^{\kappa_1b}).
\]
For the labels $P,O_\pm,U,L$, the infinite-volume factors before taking
the $m$th power are, respectively,
\[
 B(x),\qquad B(x\mp i),\qquad
 B(x+i(3/2+d))^{-1},\qquad B(x-i(3/2+d))^{-1}.
\]
Call these functions $\omega_v(x)$; the labels $M,F_{\rm pair}$ have
factor one. The explicit formula for $B$ gives $|\omega_v(x)|<1$ and
\begin{equation}\label{j:site-tail}
 \omega_v(x)=1-\Lambda_{v,\pm}e^{-\kappa_1|x|}
                 +O(e^{-2\kappa_1|x|})\quad(x\to\pm\infty),
 \qquad \Re\Lambda_{v,\pm}>0.
\end{equation}
In particular
\[
 |\omega_v(x)^m|\le e^{-cm e^{-\kappa_1|x|}},\qquad
 |\omega_v(x)^m-1|\le Cm e^{-\kappa_1|x|}.
\]
Inside $|x|\lesssim bL$, these factors suppress all labels except
$M,F_{\rm pair}$; outside they approach one. This is the logarithmic
interval over which the interior and exterior propagation rates differ.

Let $T_{m,j}$ be the one-block kernel of Proposition~\ref{j:block-kernel}
with the factor $\prod_{(t,v)\in Y}\omega_v(jb+t)^m$. Define the central
product and the exterior differences by
\begin{equation}\label{j:central-and-defects}
 \begin{split}
 \widetilde C_m&=\prod_{j=-L}^{L-1}(T_{m,j}/u),\\
 D_{+,k}&=T_{m,L+k}/r_E-E',\\
 D_{-,k}&=T_{m,-L-1-k}/r_E-E',\qquad k\ge0,
 \end{split}
\end{equation}
where products in the spatial coordinate are in increasing order. A
finite subset of exterior indices on a side has length $l_\pm$ equal to
one more than its maximum, and length zero when it is empty. Its word
$W_+$ consists of the factors on indices $0,\ldots,l_+-1$, with
$D_{+,k}$ at chosen indices and $E'$ otherwise. Define $W_-$ in the same
way in reverse order. The empty subset gives the identity operator. The
occupied index interval of a nonempty word ends at its largest chosen
index, with no factors beyond that index.

For a centered interval of $n\ge2L$ blocks containing the central range,
expand the exterior factors and rotate the trace to begin with
$C'_m=(u/r_E)^{2L}\widetilde C_m$. The term for two compatible exterior
words has the form
\begin{equation}\label{j:word-trace-form}
 \operatorname{tr}\bigl(C'_m W_+(E')^{s-l_+-l_-}W_-\bigr),
 \qquad s=n-2L.
\end{equation}
The next lemma identifies what is needed to estimate this trace: a
trace-norm bound for the central product, exponential moments of the
word lengths, and a limiting difference from a bare power trace. It
allows all peripheral nilpotent terms of $E'$.

\subsection{A coefficient lemma for signed traces}

The preceding decomposition leads to the following abstract statement.
The operator $E$ describes the exterior, $C$ contains the central
interval, and the two word families collect the perturbations at its
ends. Their lengths include the unperturbed blocks between chosen
perturbations. The lemma extracts the coefficient of each peripheral
eigenvalue from a difference of traces.

\begin{lemma}[Signed trace coefficients]\label{j:trace-extraction}
Let $E$ be a Hilbert--Schmidt operator on a complex Hilbert space with
spectral radius one. Let $\mathcal P=\{\zeta\in\sigma(E):|\zeta|=1\}$.
For each $\zeta\in\mathcal P$ let $P_\zeta$ be its Riesz projection,
$d_\zeta=\operatorname{rank}P_\zeta$, and
$J_\zeta=(E-\zeta I)P_\zeta$, with
$J_\zeta^{q_\zeta}=0$ on $\operatorname{ran}P_\zeta$.
Let $C$ be trace class. Suppose two countable families of bounded
operators $W_+,W_-$ have integer lengths $\ell(W_\pm)\ge0$ and, for
some $\epsilon>0$,
\begin{equation}\label{j:word-moment}
 \sum_{W_\pm}e^{\epsilon\ell(W_\pm)}\|W_\pm\|<\infty.
\end{equation}
For integers $s\ge0$ put
\[
 F_s=\sum_{\ell(W_+)+\ell(W_-)\le s}
 \operatorname{tr}(CW_+E^{s-\ell(W_+)-\ell(W_-)}W_-).
\]
Assume that $a\ge0$ is an integer and $R\in\mathbb C$ satisfy
$F_s-R\operatorname{tr}E^{s+a}\to0$ as $s\to\infty$.
Only sufficiently large $s$, with $s+a\ge2$, are used in this limit.
Then the coefficientwise absolutely convergent polynomial
\[
 \mathcal P_\zeta(z)=\sum_{W_+,W_-}\zeta^{-\ell}
 \operatorname{tr}\left(CW_+
 \left[\sum_{t=0}^{q_\zeta-1}
 \binom{z-\ell}{t}\zeta^{-t}J_\zeta^tP_\zeta\right]W_-\right),
 \quad \ell=\ell(W_+)+\ell(W_-),
\]
is identically $R d_\zeta\zeta^a$. In particular,
\[
 R=\frac{\mathcal P_\zeta(0)}{d_\zeta\zeta^a},\qquad
 |R|\le K\|C\|_1,
\]
where $K$ depends on $E$ and the bounds in \eqref{j:word-moment}, but
not on $C$ or $a$. More precisely it may be chosen in terms of
$\epsilon$, the two weighted word sums, the finite nilpotent orders
$q_\zeta$, and the norms $\|J_\zeta^tP_\zeta\|$.
Thus the conclusion is uniform for a family with uniform bounds on
these data. No peripheral semisimplicity is assumed.
\end{lemma}

\begin{proof}
Compactness makes $\mathcal P$ finite and nonempty. On its finite Riesz
blocks,
\[
 E^nP_\zeta=\zeta^n\sum_{t=0}^{q_\zeta-1}
 \binom nt\zeta^{-t}J_\zeta^tP_\zeta.
\]
The complementary restriction $E_<$ has spectral radius strictly less
than one. Thus the remainder in this expansion is $O(\rho^n)$ in
operator norm for some $\rho<1$, and
$\|E^n\|\le K_0(1+n)^q$ for a fixed $q$.
For $n\ge2$, the complementary trace norm is bounded by
$\|E_<\|_2^2\|E_<^{n-2}\|$. Every positive nilpotent power has trace zero
on a finite Riesz block, so
\begin{equation}\label{j:bare-trace-expansion}
 \operatorname{tr}E^n=\sum_{\zeta\in\mathcal P}
 d_\zeta\zeta^n+O(\rho_1^n),\qquad \rho_1<1.
\end{equation}

The inequality $|\operatorname{tr}(CAB)|\le\|C\|_1\|A\|\|B\|$
and the polynomial power bound prove absolute convergence of $F_s$.
A coefficient of $\binom{z-\ell}{t}$ is bounded by a constant times
$(1+\ell)^t$, so the exponential moments in \eqref{j:word-moment}
also prove coefficientwise convergence of $\mathcal P_\zeta$.
For pairs with $\ell\le s/2$, replacing the middle power by its
peripheral expansion costs $O(\rho^{s/2})\|C\|_1$. For the remaining
pairs use the polynomial bound and the exponential tail of the word
sums. Their total is at most a polynomial in $s$ times
$e^{-\epsilon s/2}\|C\|_1$. The same estimate permits extending the
peripheral sum to all word pairs, using the polynomial growth of the
generalized binomial coefficients. Hence
\[
 F_s=\sum_{\zeta\in\mathcal P}\zeta^s\mathcal P_\zeta(s)+o(1).
\]
Together with the assumed difference of traces and
\eqref{j:bare-trace-expansion}, this gives
\[
 \sum_{\zeta\in\mathcal P}\zeta^s
 \big(\mathcal P_\zeta(s)-Rd_\zeta\zeta^a\big)\longrightarrow0.
\]
Distinct unit-modulus exponentials cannot cancel a nonzero polynomial
coefficient in this way. Divide by the largest occurring power $s^d$,
multiply by $\overline\eta^{s}$ for one $\eta\in\mathcal P$, and take
Ces\`aro means. The geometric-series formula makes every frequency
except $\eta$ vanish, leaving its leading coefficient. That coefficient
is zero; descend through the degrees. This proves the polynomial
identity. Evaluation at zero is bounded by
\[
 K_1\|C\|_1\sum_{W_+,W_-}
 \|W_+\|\|W_-\|(1+\ell)^{q_\zeta-1},
\]
which is finite by \eqref{j:word-moment}. Since $d_\zeta\ge1$ and
$|\zeta^a|=1$, the asserted bound follows.
\end{proof}


\subsection{Bounds for the central product and the words}

Fix $b,d$ and a compact parameter neighborhood $D$ inside
$|\Re\theta|<\pi/2$. The gauged kernel bound and \eqref{j:site-tail} give
\begin{equation}\label{j:defect-bounds}
 \|D_{\pm,k}\|_2\le Ce^{-\kappa_1bk}.
\end{equation}
For $-L\le j<L$, the Hilbert--Schmidt norm of
$T_{m,j}/u-I_0/u$ is at most
$C\exp(-c e^{\kappa_1bk})$, where $k$ is the number of full blocks
between $j$ and the nearer edge. These estimates follow by telescoping
the products of the specified scalar multipliers. A telescoping term
costs at most a particle count times the displayed scalar difference;
the quadratic gauge bound absorbs that count. In the central estimate,
any configuration omitted by $I_0$ contains at least one suppressed
particle.

Expand the central differences against the bounded powers
\eqref{j:interior-powers}. Their norms are summable from both ends.
Keeping a Hilbert--Schmidt factor at each endpoint gives
\begin{equation}\label{j:central-trace-bound}
 \|\widetilde C_m\|_1\le C.
\end{equation}
Likewise the peripheral projections and nilpotent powers of $E'$ are
uniformly bounded on $D$ by separating contours, with orders bounded by
their fixed finite ranks. Thus $\|(E')^k\|\le A(1+k)^{r_J}$ there, for a
fixed nonnegative integer $r_J$. For a nonempty exterior word, pair each
intervening power with the farther
chosen perturbation at index $k$. There is no unrestricted gap, so
\[
 \|W_\pm\|\le
 \prod_{k\text{ chosen}}A_1(1+k)^{r_J} e^{-\kappa_1bk}.
\]
Its length is at most the sum of $k+1$ over its chosen indices. Hence for
$0<\epsilon<\kappa_1b$,
\begin{equation}\label{j:exterior-word-bound}
 \sum_W e^{\epsilon l(W)}\|W\|
 \le\prod_{k\ge0}
 \left(1+A_1(1+k)^{r_J} e^{-\kappa_1bk+\epsilon(k+1)}\right)<\infty,
\end{equation}
uniformly in $m$. Equations~\eqref{j:central-trace-bound} and
\eqref{j:exterior-word-bound} give the uniform central and word bounds
required by Lemma~\ref{j:trace-extraction}.

We also need the dependence on the position of each edge within a block.
At the right edge, with $s=x-bL$, replace the suppressed powers by
$\exp(-a\Lambda_{v,+}e^{-\kappa_1s})$. At the left edge, with $s=x+bL$,
use $\exp(-a\Lambda_{v,-}e^{\kappa_1s})$. These profiles define kernels
analytic in $a$ near the closed interval $[1,e^{\kappa_1b}]$. Choose this
neighborhood so that $\Re(a\Lambda_{v,\pm})$ remains positive for every
suppressed label; the same domination and summable bounds then apply.
The exterior differences in
\eqref{j:central-and-defects} converge to their profile differences at
$a=a_m$ with error $Ce^{-cL}e^{-ck}$. In the first
$\lfloor L/4\rfloor$ central blocks from each edge the analogous
normalized error is $Ce^{-cL}$. Indeed the error from the two-term
expansion of $\log\omega_v$, after multiplication by $m$, is
$O(me^{-2\kappa_1|x|})$ in these regions. Deeper central blocks differ
from $I_0/u$ by a superexponentially small amount.

It follows that there is an analytic trace-class family $C(a,\theta)$
such that
\begin{equation}\label{j:central-profile}
 \widetilde C_m=C(a_m,\theta)+O(e^{-cL})
 \quad\hbox{in trace norm}.
\end{equation}
Construct it by taking the ordered left interior profile from its edge
inwards, then $P_I$, then the right interior profile from inside to its
edge, and taking limits in the two lengths. Inserting an additional block
next to $P_I$ changes the product by a summable superexponential error,
since $P_I$ is fixed by $I_0/u$ and all partial products are bounded.
For comparison with $\widetilde C_m$, the middle portion after the first
interior blocks from both edges is $P_I+O(e^{-cL})$ in operator norm by
\eqref{j:interior-powers} and the same expansion. The bracketing products
have bounded Hilbert--Schmidt norms. Replacing their edge portions by the
profile kernels then gives \eqref{j:central-profile} by telescoping.

\subsection{The finite-site limit and coefficient extraction}

For fixed $m$, define the normalized physical quantity
\begin{equation}\label{j:normalized-cylinder}
 R_m(\theta)=e^{-m^2\kappa_\infty}\frac{N(X)^2}{P_0}
             G_{A,B}(h+h^{-1}),\qquad
 \kappa_\infty=\lim_{T\to\infty}\kappa_T,
\end{equation}
with homogeneous $X$ and $h=e^{-2i\theta}$. The fixed-site contour limit,
the gas identity, and the $m=0$ trace in \eqref{j:empty-block-traces}
identify $R_m$ with the $n\to\infty$ limit of the ratio of the trace with
periodized site factors to $\operatorname{tr}E_0^n$. Use blocks
$-\lfloor n/2\rfloor\le j<\lceil n/2\rceil$.

In the numerator normalized by $r_E^n$, the periodized site factors can
be replaced by the actual single-profile factors $T_{m,j}$, not by their
edge limits, with an error tending to zero. At fixed $m$ each block
difference costs $O_m(e^{-cn})$ in the kernel estimate. To control its
surrounding products, expand the unperiodized exterior differences by
subsets in increasing or decreasing exterior index. Pair each power gap
between successive perturbations with the farther index; the polynomial
power bound for $E'$ is then absorbed by the exponential decay in
\eqref{j:defect-bounds}. Only one unrestricted power gap remains, giving
a polynomial bound in $n$ for any exterior half-segment. A whole
subinterval has at most two such portions and a fixed central part.
Expanding the additional periodization differences retains these
polynomial bounds, while their exponential smallness makes the total
error tend to zero. The same argument applies to traces by keeping two
Hilbert--Schmidt factors.

The bare normalized traces $\operatorname{tr}(E')^n$ are bounded: their
peripheral expansions contain only multiplicity times eigenvalue powers,
since positive nilpotent powers have trace zero. Multiplying the known
ratio limit by these bounded traces and using the replacement just proved
therefore gives
\begin{equation}\label{j:fixed-site-trace-difference}
 \operatorname{tr}\prod_j(T_{m,j}/r_E)
       -R_m\operatorname{tr}(E')^n\longrightarrow0
\end{equation}
at fixed $m$. No lower bound on the bare trace is used.

We now apply Lemma~\ref{j:trace-extraction} to
\eqref{j:word-trace-form}. Each actual exterior half has length
$s/2+O(1)$, rather than allowing every pair with $l_++l_-\le s$.
Extending to that latter set adds only pairs with total word length at
least $s/2-O(1)$. By \eqref{j:exterior-word-bound} and the polynomial
power bound, their contribution is $O_m((1+s)^r e^{-cs})$. It is
negligible at fixed $m$. Thus \eqref{j:fixed-site-trace-difference}
is the lemma's difference hypothesis with
\[
 E=E',\qquad C=C'_m=(u/r_E)^{2L}\widetilde C_m,
 \qquad a=2L.
\]

For a peripheral eigenvalue $\zeta$ of $E'$, write $P_\zeta$ for its
Riesz projection, $d_\zeta=\operatorname{rank}P_\zeta$, and
$J_\zeta=(E'-\zeta)P_\zeta$. The coefficient polynomial supplied by the
lemma is, in the present notation,
\begin{equation}\label{j:cylinder-coefficient}
 \mathcal P_{\zeta,m}(s)=
 \sum_{W_+,W_-}\zeta^{-l_+-l_-}
 \operatorname{tr}\!\left(
 C'_mW_+\left[\sum_{t\ge0}
 \binom{s-l_+-l_-}{t}\zeta^{-t}J_\zeta^tP_\zeta\right]W_-\right).
\end{equation}
The nilpotent sum is finite. The word sum is coefficientwise absolutely
convergent by \eqref{j:exterior-word-bound}. The lemma gives directly
\begin{equation}\label{j:balanced-trace-estimate}
 R_m=\frac{\mathcal P_{\zeta,m}(0)}{d_\zeta\zeta^{2L}},
 \qquad |R_m|\le C|u/r_E|^{2L}.
\end{equation}
The constant is uniform in $m$ on the fixed parameter neighborhood:
the word moments, central trace norm, and peripheral nilpotent bounds are
uniform, and $|\zeta^{2L}|=1$ makes the changing offset harmless.

The analytic dependence on the edge position is also retained. Replace
$\widetilde C_m$ in \eqref{j:cylinder-coefficient} by $C(a,\theta)$ and
each exterior difference by its profile limit. Their exponentially small
weighted errors telescope in the absolutely summable word series. Hence,
locally uniformly in $|\Re\theta|<\pi/2$,
\begin{equation}\label{j:balanced-trace-profile}
 R_m(u/r_E)^{-2L}d_\zeta\zeta^{2L}
 =\mathcal F_\zeta(a_m,\theta)+O(e^{-cL}),
\end{equation}
where $\mathcal F_\zeta$ is analytic in $a$ near
$[1,e^{\kappa_1b}]$ and locally analytic in $\theta$. Fixed separating
contours give the stated uniform spectral bounds. Equations
\eqref{j:balanced-trace-estimate}--\eqref{j:balanced-trace-profile} are
the operator outputs used after the remaining homogeneous scalar
normalization.
